\documentclass[11pt]{article}

\usepackage[utf8]{inputenc}
\usepackage[T1]{fontenc}
\usepackage[margin=1in]{geometry}
\usepackage{array}
\usepackage{longtable}
\usepackage{booktabs}

\title{Checklist de Cumplimiento -- Laboratorio 3}
\author{}
\date{}

\begin{document}

\maketitle

\noindent Este archivo no es para entregar como informe. Su funci\'on es verificar, de forma trazable, que el archivo principal \texttt{informe\_lab3\_estado\_ieee.tex} cubre cada criterio de la r\'ubrica y cada pregunta de las gu\'ias.

\renewcommand{\arraystretch}{1.15}
\begin{longtable}{p{2.2cm} p{8.7cm} p{4.0cm}}
\toprule
Bloque & Criterio de la r\'ubrica o de la gu\'ia & Ubicaci\'on en el informe principal \\
\midrule
\endfirsthead
\toprule
Bloque & Criterio de la r\'ubrica o de la gu\'ia & Ubicaci\'on en el informe principal \\
\midrule
\endhead
\midrule
\multicolumn{3}{r}{Contin\'ua en la siguiente p\'agina}
\endfoot
\bottomrule
\endlastfoot

3.1 sim & Mostrar ecuaciones linealizadas en el punto de operaci\'on & Secci\'on ``Resultados'', subsecci\'on ``Pr\'actica 3.1: resultados de simulaci\'on'', bloque ``Ecuaciones linealizadas en el punto de operaci\'on'' \\
3.1 sim & Identificar matrices $A$, $B$, $C$ y $D$ & Misma subsecci\'on, bloque ``Representaci\'on en espacio de estados del sistema linealizado''; se reporta la implementaci\'on final con dos salidas y tambi\'en la salida m\'inima $C_{\alpha}$ \\
3.1 sim & Analizar estabilidad a partir de los polos & Secci\'on ``An\'alisis'', subsecci\'on ``Pr\'actica 3.1''; los polos se reportan en Resultados \\
3.1 sim & Hallar matriz de controlabilidad y concluir si es controlable & Se reporta en Resultados y se interpreta en An\'alisis \\
3.1 sim & Hallar matriz de observabilidad y concluir si es observable & Se reporta en Resultados y se interpreta en An\'alisis \\
3.1 sim & Hallar polos dominantes con $\zeta=0.7$, $\omega_n=4$ & Secci\'on ``Resultados'', subsecci\'on ``Pr\'actica 3.1: resultados de simulaci\'on'' \\
3.1 sim & Hallar polinomio caracter\'istico de orden 4 con polos adicionales en $-30$ y $-40$ & Misma subsecci\'on \\
3.1 sim & Hallar vector $K$ del controlador & Misma subsecci\'on \\
3.1 sim & Mostrar diagrama de bloques con realimentaci\'on de estados & Figura etiquetada como \texttt{fig:p31\_bloques} en el informe principal \\
3.1 sim & Mostrar gr\'afica de simulaci\'on de referencia y salida de $x_1$ & Figura etiquetada como \texttt{fig:p31\_sim}, panel (a) \\
3.1 sim & Mostrar gr\'afica de simulaci\'on de referencia y salida de $x_2$ & Figura etiquetada como \texttt{fig:p31\_sim}, panel (b) \\
3.1 sim & Mostrar gr\'afica de simulaci\'on de la se\~nal de control & Figura etiquetada como \texttt{fig:p31\_sim}, panel (c) \\

3.1 exp & Mostrar gr\'afica experimental de referencia y salida de $x_1$ & Figura etiquetada como \texttt{fig:p31\_exp}, panel (a) \\
3.1 exp & Mostrar gr\'afica experimental de referencia y salida de $x_2$ & Figura etiquetada como \texttt{fig:p31\_exp}, panel (b) \\
3.1 exp & Mostrar gr\'afica experimental de la se\~nal de control & Figura etiquetada como \texttt{fig:p31\_exp}, panel (c) \\
3.1 exp & Comparar simulaci\'on y experimento & Secci\'on ``An\'alisis'', subsecci\'on ``Pr\'actica 3.1'' \\
3.1 preg & Responder implicaciones de controlabilidad y observabilidad en realimentaci\'on de estados & Secci\'on ``An\'alisis'', subsecci\'on ``Pr\'actica 3.1'' \\
3.1 preg & Responder qu\'e pasa si la condici\'on inicial est\'a lejos del equilibrio & Secci\'on ``An\'alisis'', subsecci\'on ``Pr\'actica 3.1'' \\

3.2 sim & Mostrar ecuaciones linealizadas en el punto de operaci\'on & Secci\'on ``Resultados'', subsecci\'on ``Pr\'actica 3.2: resultados de simulaci\'on'', bloque ``Ecuaciones linealizadas en el punto de operaci\'on'' \\
3.2 sim & Reportar matrices $A$, $B$ y $C$ & Misma subsecci\'on, bloque ``Representaci\'on en espacio de estados del sistema linealizado'' \\
3.2 sim & Demostrar controlabilidad & Se reporta en Resultados y se interpreta en An\'alisis \\
3.2 sim & Demostrar observabilidad & Se reporta en Resultados y se interpreta en An\'alisis \\
3.2 sim & Reportar polinomio caracter\'istico $(s+100)(s^2+28s+400)$ & Secci\'on ``Resultados'', subsecci\'on ``Pr\'actica 3.2: resultados de simulaci\'on'' \\
3.2 sim & Reportar vector $K$ del controlador & Misma subsecci\'on \\
3.2 sim & Mostrar diagrama del controlador en espacio de estados & Figura etiquetada como \texttt{fig:p32\_bloques}, panel (a) \\
3.2 sim & Mostrar diagrama con ganancia $K_g$ & Figura etiquetada como \texttt{fig:p32\_bloques}, panel (b) \\
3.2 sim & Mostrar diagrama con control integral & Figura etiquetada como \texttt{fig:p32\_bloques}, panel (c) \\
3.2 sim & Mostrar gr\'afica de salida con referencia sin $K_g$ & Figura etiquetada como \texttt{fig:p32\_sim}, panel (a) \\
3.2 sim & Mostrar gr\'afica de salida con referencia con $K_g=-705$ & Figura etiquetada como \texttt{fig:p32\_sim}, panel (b) \\
3.2 sim & Mostrar gr\'afica de salida con referencia con control integral & Figura etiquetada como \texttt{fig:p32\_sim}, panel (c) \\
3.2 sim & Analizar el caso sin $K_g$, con $K_g$ y con integral & Secci\'on ``An\'alisis'', subsecci\'on ``Pr\'actica 3.2'' \\

3.2 exp & Mostrar gr\'afica experimental con referencia y salida & Figura etiquetada como \texttt{fig:p32\_exp} \\
3.2 exp & Analizar la gr\'afica experimental y el comportamiento & Secci\'on ``An\'alisis'', subsecci\'on ``Pr\'actica 3.2'' \\
3.2 exp & Comparar experimento con simulaci\'on & Secci\'on ``An\'alisis'', subsecci\'on ``Pr\'actica 3.2'' \\
3.2 preg & Responder ventajas de la representaci\'on de estados para implementar un controlador & Secci\'on ``An\'alisis'', subsecci\'on ``Pr\'actica 3.2'' \\

Formato & Mantener las secciones principales de introducci\'on, resultados, an\'alisis y conclusiones & Estructura del archivo principal \\
Formato & Incluir las im\'agenes finales y referirlas desde el an\'alisis & Figuras \texttt{fig:p31\_bloques}, \texttt{fig:p31\_sim}, \texttt{fig:p31\_exp}, \texttt{fig:p32\_bloques}, \texttt{fig:p32\_sim}, \texttt{fig:p32\_exp} \\
\end{longtable}

\end{document}
